\documentclass[10pt,a4paper]{amsart}
\usepackage[margin=19mm]{geometry}
\usepackage{amsmath,amssymb,mathtools}
\usepackage[scr=rsfs,cal=euler]{mathalfa}
\usepackage{xfrac}
\usepackage[unicode,hidelinks,bookmarks=false]{hyperref}
\newcommand{\ie}{i.e.\ }
\newcommand{\cf}{cf.\ }
\newcommand{\resp}{respectively\ }
\newcommand{\Subsection}{\subsection}
\providecommand{\nobreakdash}{\nobreak\hbox{-}}
\setlength{\emergencystretch}{2em}
\multlinegap0pt
\usepackage{bm}
\usepackage{bbm}
\usepackage{dsfont}
\usepackage{xspace}
\usepackage{cancel}
\usepackage{mathtools}
\usepackage{amsthm}
\makeatletter
\@ifundefined{theorem}{\newtheorem{theorem}{Theorem}}{}
\@ifundefined{lemma}{\newtheorem{lemma}[theorem]{Lemma}}{}
\makeatother
\DeclareMathOperator{\wmaj}{wmaj}
\newcommand{\posi}{{\gamma}}
\newcommand{\pos}{{\boldsymbol\gamma}}
\newcommand{\rem}{{\boldsymbol\rho}}
\title[A generalization of conjugation of integer partitions]{A generalization of conjugation of integer partitions}
\author{Seamus Albion, Theresia Eisenk\"olbl, Ilse Fischer, Moritz Gangl, Hans H\"ongesberg, Christian Krattenthaler, Martin Rubey}
\date{Source: arXiv:2407.16043v2 (CC BY 4.0)}
\begin{document}
\maketitle

\noindent\emph{Source and licence.} A generalization of conjugation of integer partitions. Version arXiv:2407.16043v2 (CC BY 4.0), distributed under the Creative Commons Attribution 4.0 International licence, as recorded in the arXiv licence metadata for this exact version and shown on the abstract page https://arxiv.org/abs/2407.16043v2. Full rights notice: licenses/arxiv-2407.16043-CC-BY-4.0.txt.\par\smallskip

\noindent\textit{Editorial scope.} Counted unit: the Lemma whose source label is \texttt{maj} in the article (source0.tex lines 958--966, source label \texttt{maj}), delivered together with the article's own complete original proof of it (source0.tex lines 968--1025, one \texttt{proof} environment of 58 lines). No mathematics is written, repaired or re-derived: statement and proof are the article's own text line for line. Required context retained uncounted: none. Every definition, display and label of those lines is retained unchanged. Omitted from this package and not redistributed: 1--957, 967--967, 1026--1724. Nothing inside a retained excerpt is omitted or altered: every retained body is delivered exactly as the article's lines are written, so no omission receipt of any kind accompanies this package. Citations: the retained excerpts carry no citation command at all, so no bibliography entry of the article is transcribed and no reference is redistributed. Reference adaptation: none was needed and none was made; every reference inside the delivered unit resolves inside it, so no unresolved reference or citation is produced. Printed slips kept literal: no printed word, symbol, hypothesis or number of the counted statement or of its proof was altered. Layout and class: the article's own class is \texttt{amsart} with packages including inputenc, amsmath, amsthm, amssymb, url, mathdots, verbatim, geometry, xcolor, tikz, ytableau, hyperref, MnSymbol, graphicx; the wrapper is the lane's pinned \texttt{amsart} editorial wrapper at 10pt on A4, which restates the theorem-like environments the retained text needs and adds the article's own macro definitions that the retained text needs (\texttt{\textbackslash pos}, \texttt{\textbackslash posi}, \texttt{\textbackslash rem}, \texttt{\textbackslash wmaj}), with extra wrapper packages (bm, bbm, dsfont, xspace, cancel, mathtools). The delivered numbering is the unit's own. Assets: no retained span contains a figure environment, a graphics-inclusion command, a PDF-page inclusion, a table or a TikZ picture; no image file is copied into this package, the package declares no asset dependency and no image file is redistributed with it. Licence: version arXiv:2407.16043v2 of this article is distributed under the Creative Commons Attribution 4.0 International licence, as recorded in the arXiv licence metadata for this exact version and shown on the abstract page https://arxiv.org/abs/2407.16043v2. A full rights notice is delivered at licenses/arxiv-2407.16043-CC-BY-4.0.txt. Characters: every retained excerpt is pure ASCII, so the delivered engine, XeLaTeX with its default Latin Modern text font, reports no missing character.
\begin{lemma}
\label{maj}
Let $\rem$ be a vector of integers between $1$ and~$s-1$ of length~$m$. Then, for fixed $\posi_m$, we have
$$
\sum_{1 \le \posi_1 < \posi_2 < \dots < \posi_m} Q^{{\mathbf d}(\rem,\pos) \cdot ({\pos- \mathbf 1})} =
Q^{-\wmaj(\rem)} \sum_{1 \le \posi_1 < \posi_2 < \dots < \posi_m} Q^{|\pos|-m},
$$
where $\pos=(\posi_1,\dots,\posi_m)$.
\end{lemma}

\begin{proof}
We need the following generalization of the weak major index: for $k$ with $1 \le k < m$,
we define
$$
\wmaj_k(\rem)= \underset{j\le k}{\sum_{j:\rho_j \ge \rho_{j+1}}} j.
$$
This simply is the weak major index of the tuple $\rem$ cut off after
the {$(k+1)$-st entry.} Note that $\wmaj_{m-1}=\wmaj$ for sequences of
length~$m$. 

The proof is by induction on~$m$. For the start of the
induction we note that for $m=1$ the statement is obvious.

In the following arguments, the reader should always keep in mind
that, by assumption, $\posi_m$ is fixed throughout.
In particular, if we write a sum 
$\sum_{1 \le \posi_1 < \posi_2 < \dots < \posi_m} \dots$, then the
sum runs over the $\posi_i$ with $i<m$, while $\posi_m$ is fixed. 

Now, by the induction hypothesis, we may assume
\begin{multline*}
\sum_{1 \le \posi_1 < \posi_2 < \dots < \posi_m} Q^{{\mathbf d}(\rem, \pos) \cdot ({\pos- \mathbf 1})}
=\sum_{\posi_{m-1}<\posi_{m}} Q^{d_m(\gamma_m -1)} \sum_{1 \le \posi_1 < \posi_2 < \dots < \posi_{m-1}} Q^{\sum_{j=1}^{m-1} d_j(\posi_j-1)}\\ 
=Q^{-\wmaj_{m-2}(\rem)}
\sum_{1 \le \posi_1 < \posi_2 < \dots < \posi_m} Q^{\sum_{j=1}^{m-1} (\posi_j-1) + d_{m} (\posi_{m}-1)}.
\end{multline*}
If $\rho_{m-1} < \rho_m$, then $d_m=1$ and
$\wmaj(\rem)=\wmaj_{m-2}(\rem)$, and the assertion follows in this
case. If, on the other hand, we have $\rho_{m-1} \ge \rho_m$, then
$\wmaj(\rem)=\wmaj_{m-2}(\rem)+m-1$, and, by the definition of $d_m$, we
have
\begin{multline}
\label{left}
\sum_{1 \le \posi_1 < \posi_2 < \dots < \posi_{m-1} < \posi_m} Q^{\sum_{j=1}^{m-1} (\posi_j-1) + d_{m} (\posi_{m}-1)} \\
= \sum_{1 \le \posi_1 < \posi_2 < \dots < \posi_{m-1} < \posi_m-1} Q^{|\pos|-m}
+ \sum_{ 1 \le \posi_1 < \posi_2 < \dots < \posi_{m-2} < \posi_m-1} Q^{\sum_{j=1}^{m-2} (\posi_j-1)+\posi_{m}-2},
\end{multline}
where the first sum of the right-hand side corresponds to the case $\posi_m > \posi_{m-1}+1$ and the second sum corresponds to the case $\posi_m = \posi_{m-1}+1$.

We need to show that this is equal to
$$
Q^{-m+1} \sum_{1 \le \posi_1 < \posi_2 < \dots < \posi_m} Q^{|\pos|-m}.
$$
We provide a combinatorial proof.
First note that the first term in the second line of \eqref{left} can be transformed as follows:
$$
\sum_{1 \le \posi_1 < \posi_2 < \dots < \posi_{m-1} < \posi_m-1} Q^{|\pos|-m} =
Q^{-m+1} \sum_{2 \le \posi_1 < \posi_2 < \dots < \posi_{m-1} < \posi_m} Q^{|\pos|-m}.
$$
Here we have used the transformation $\posi_i \to \posi_i-1$ for $i \in \{1,2,\dots,m-1\}$. The second term
in the second line of \eqref{left} is
$$
\sum_{1 \le \posi_1 < \posi_2 < \dots < \posi_{m-2} < \posi_m-1} Q^{\sum_{j=1}^{m-2} (\posi_j-1)+\posi_{m}-2} =
Q^{-m+1} \sum_{1=\posi_1 < \posi_2 < \dots < \posi_{m-1} < \posi_m} Q^{|\pos|-m},
$$
where we have used the transformation $\posi_i \to \posi_{i+1}-1$ for $i \in \{1,2,\dots,m-2\}$ and have set $\posi_1=1$.
This completes the proof.
\end{proof}

\end{document}
